% RM 058 — GeMeA Status: Revised Roadmap, Mocho v1.1, Paper
% Presenter: Mary Ann Tan, FIZ-KDAI
% Date: 29 April 2026
\documentclass[xcolor,aspectratio=169,compress]{beamer}
\usetheme{Boadilla}
\definecolor{maincolor}{rgb}{.38,.06,.80}
\usecolortheme[named=maincolor]{structure}
\usepackage{sansmath}
\usepackage{svg}
\usepackage{tikzsymbols}
\usepackage{array}
\usepackage{booktabs}
\usepackage{colortbl}
\usepackage{pifont}

\newcommand{\cmark}{\textcolor{green!60!black}{\ding{51}}}
\newcommand{\xmark}{\textcolor{red}{\ding{55}}}
\newcommand{\done}{\textcolor{green!60!black}{Done}}
\newcommand{\current}{\textcolor{orange}{\textbf{Current}}}
\newcommand{\blocked}{\textcolor{red}{Blocked}}
\newcommand{\deferred}{\textcolor{gray}{Deferred}}

\makeatother
\setbeamertemplate{footline}
{
  \leavevmode%
  \hbox{%
  \begin{beamercolorbox}[wd=.30\paperwidth,ht=2.25ex,dp=1ex,center]{author in head/foot}%
    \usebeamerfont{author in head/foot}\insertshortauthor{}
    (\insertshortdate)
  \end{beamercolorbox}%
  \begin{beamercolorbox}[wd=.55\paperwidth,ht=2.25ex,dp=1ex,center]{title in head/foot}%
    \usebeamerfont{title in head/foot}\insertshorttitle\hspace*{3em}
  \end{beamercolorbox}%
  \begin{beamercolorbox}[wd=.15\paperwidth,ht=2.25ex,dp=1ex,right]{date in head/foot}%
   \insertframenumber{} / \inserttotalframenumber\hspace*{1ex}
  \end{beamercolorbox}}%
  \vskip0pt%
}
\makeatletter

\setbeamertemplate{headline}[default]
\setbeamertemplate{navigation symbols}{}
\mode<beamer>{\setbeamertemplate{blocks}[rounded][shadow=false]}
\setbeamercovered{transparent}
\useoutertheme[subsection=false]{miniframes}

\title[RM 058 --- GeMeA Status]{GeMeA Status Update}

\author[M.A.\ Tan]{Mary Ann Tan}

\institute[FIZ-KDAI]{FIZ Karlsruhe -- Leibniz Institute for Information Infrastructure, KDAI}

\date[RM 058, 29.04.2026]{Research Meeting 058, {\sansmath 29$^{\rm th}$} April 2026}

% \logo{\includesvg[width=.2\textwidth]{../beamer-bodilla-template/Logo/FIZ_Logo_en.svg}\vspace{215pt}}

% % % % % %
\begin{document}

\frame[plain]{\titlepage}

% ---- OUTLINE ----
\begin{frame}{Outline}
  \begin{enumerate}
    \item Revised Roadmap
    \item Mocho v1.1
    \item NER Annotation
    \item Status \& Deliverables
  \end{enumerate}
\end{frame}

% =================================================================
\section{Revised Roadmap}
% source: gemea/notes/project/roadmap-revised.md
% =================================================================

% ---- 1.1 What Changed ----
% source: roadmap-revised.md §1 Scope
\begin{frame}{Revised Roadmap --- What Changed}
  {\tiny\url{https://github.com/anntanp/gemea/blob/develop/notes/project/roadmap-revised.md}}
  \small
  \begin{columns}[t]
    \begin{column}{.48\textwidth}
      \textbf{Original scope (dropped / deferred):}
      \begin{itemize}\setlength{\itemsep}{2pt}
        \item[$\times$] Elasticsearch + faceted search $\to$ v2
        \item[$\times$] Next.js KG browser frontend $\to$ v2
        \item[$\times$] FastAPI / GraphQL API $\to$ v2
        \item[$\times$] Full 65M-object corpus $\to$ v2
        \item[$\times$] GND Works linking (NER) $\to$ described only; not implemented
      \end{itemize}
    \end{column}
    \begin{column}{.50\textwidth}
      \textbf{v1 scope (what we are claiming):}
      \begin{itemize}\setlength{\itemsep}{2pt}
        \item[$\checkmark$] \textbf{27M} DDB objects in QLever, public SPARQL endpoint
        \item[$\checkmark$] \textbf{Mocho PoC} in QLever (representative subset)
        \item[$\checkmark$] \textbf{SHMARQL} --- lightweight Linked Data browser over QLever
        \item[$\checkmark$] \textbf{MCP/MCPO} --- agentic access layer on QLever
        \item[$\circ$] GND agent linking: described as pipeline component; deferred to v1.1
      \end{itemize}
    \end{column}
  \end{columns}
  \smallskip
  \begin{block}{Core claim}
    \footnotesize First openly browsable, SPARQL-accessible, and MCP-accessible KG over DDB
    (27M objects) with mocho-based ontology alignment PoC; testbed for LLM + KG research
    at cultural heritage scale.
  \end{block}
\end{frame}

% ---- 1.1b Revised Contribution Framing ----
% source: gemea/notes/paper/plan-abstract.md v3
%         gemea/notes/paper/paper-framing.md §Suggested framing adjustment
%         gemea/notes/project/roadmap-revised.md §2 Framing + §8 What Determines Acceptance
\begin{frame}{Contribution}
  \small
  \begin{itemize}\setlength{\itemsep}{4pt}
    \item First openly browsable, \textbf{SPARQL-accessible, and MCP-accessible} KG over DDB (27M objects)
    \item Lightweight Linked Data browser (SHMARQL) --- zero custom frontend
    \item mocho alignment PoC: EDM $\to$ \textbf{domain- and modality-specific ontologies}
          (RDA, RiC-O, MO/ACO, VRA, EBUCore, DoCO\ldots) on a quantified, scoped subset
  \end{itemize}
  \smallskip
  \footnotesize
  \begin{tabular}{lp{9cm}}
    \toprule
    1. Novel access layer & MCP interface on a CH KG --- no prior work has this \\
    2. Multi-domain alignment PoC & Data-driven EDM $\to$ sector/modality ontologies via mocho; quantified + honestly scoped \\
    3. Testbed / Dataset for Future Research & enables studies on ontology evaluation, agentic KG (MCP), RML mapping, LLM-assisted KG enrichment \\
    \bottomrule
  \end{tabular}
\end{frame}

% ---- 1.2 Roadmap Overview ----
% source: roadmap-revised.md §1, §3; rm057 roadmap table
\begin{frame}{Roadmap Overview --- v1 vs.\ v2}
  \footnotesize
  \begin{tabular}{cl>{\raggedright\arraybackslash}p{6.5cm}l}
    \toprule
    Phase & Name & v1 scope & Status \\
    \midrule
    0  & Data Acquisition   & DDB JSON-LD ingest; \texttt{langid} filtering; sector ID files & \done \\
    1  & Ingest             & QLever index (27M objects); mocho PoC; named graphs; SHMARQL  & \current \\
    1b & GND Enrichment     & \texttt{link\_gnd\_agents.py}: described, not implemented      & \blocked \\
    2  & API                & MCP/MCPO on QLever (replaces FastAPI for v1)                   & \current \\
    3  & Frontend           & SHMARQL (replaces Next.js for v1); full browser $\to$ v2       & \current \\
    4  & DevOps             & Docker Compose; QLever + SHMARQL + MCPO; CI/CD                 & \current \\
    \bottomrule
  \end{tabular}
\end{frame}

% ---- 1.3 Current Status ----
% source: roadmap-revised.md §3
\begin{frame}{Current Status}
  \small
  \begin{tabular}{ll>{\raggedright\arraybackslash}p{5.2cm}}
    \toprule
    Component & Status & Note \\
    \midrule
    QLever endpoint (115k objects, 8M triples)
      & \textcolor{orange}{\textbf{Local only}} $\triangle$
      & not yet public \\
    Mocho PoC in QLever
      & \textcolor{green!60!black}{Running} \cmark
      & subset; classes + triples TBD \\
    SHMARQL over QLever
      & \textcolor{green!60!black}{Pattern proven} \cmark
      & goethe-faust; adapt for GeMeA \\
    MCP/MCPO on QLever
      & \textcolor{green!60!black}{Pattern proven} \cmark
      & needs tool inventory \\
    GND Linking (Werk + Agent)
      & \textcolor{red}{Blocked}
      & scope: Monograph, Multi-volume, Manuscript only;
        \footnotesize chapters/sections/articles removed per Q2--Q3 2026
        feed restructure (Maria, FIZ-IR) \\
    Persistent URI (w3id / Zenodo)
      & \textcolor{red}{\textbf{Needed before 2 May}}
      & mandatory for abstract \\
    \bottomrule
  \end{tabular}
  \smallskip

  % source: roadmap-revised.md §3.1 ID corpus by sector
  \par\footnotesize\textbf{teach03:} 27.4M objects ---
  Library 18.6M, Archive 3.6M, Museum 2.0M.
\end{frame}

% ---- 1.3b DDB Feed Restructure: Risk to Use Case ----
% source: mocho-gatherer-adr.md §4.2 R6; Maria (FIZ-IR), April 2026
\begin{frame}[shrink=10]{DDB Feed Restructure --- Risk to Collection Browse Use Case}
  \small
  % source: mocho-gatherer-adr.md §4.2 R6 — Work-level top result, item-level drill-down
  \textbf{Planned use case (R6):} users browse at two levels ---
  \begin{itemize}\setlength{\itemsep}{2pt}
    \item \textbf{Top level:} primary objects / topmost document component
          (Monograph, Multi-volume work, Manuscript)
    \item \textbf{Drill-down:} sub-components mapped via \texttt{htype} to document-structure classes ---
          e.g.\ \texttt{htype\_018} (Kapitel) $\to$ \texttt{doco:Chapter},
          \texttt{htype\_006} (Aufsatz) $\to$ \texttt{doco:Section}
          (see \texttt{goethe-faust/output/lookup\_htype\_doco\_rico.csv})
  \end{itemize}

  \medskip
  \textbf{DDB planned change (Maria, FIZ-IR, Q2--Q3 2026):}
  \begin{block}{}
    \footnotesize
    Objects with \texttt{htype} = Kapitel, Abschnitt, Aufsatz, etc.\ will no longer be ingested
    after the Splitter is removed and data is reloaded.
    Only the associated Monographien will remain in the DDB.
    The change will be applied incrementally --- dataset by dataset on reload.
  \end{block}

  \begin{alertblock}{Risk}
    \small Sub-component records (chapters, sections, articles) will disappear from the DDB feed.
    The drill-down layer of the browse use case loses its primary source of
    fine-grained document-structure objects.
    GeMeA must decide whether to: (a) preserve sub-components from earlier snapshots,
    (b) scope the browse use case to top-level objects only, or
    (c) derive sub-components from full-text / structural metadata.
  \end{alertblock}
\end{frame}

% ---- 1.4 PoC Pipeline ----
% source: goethe-faust/notes/openwebui-ollama-setup.md
\begin{frame}{PoC Pipeline --- OpenWebUI + Ollama + QLever}
  \small
  {\tiny\url{https://github.com/anntanp/goethe-faust/blob/main/notes/openwebui-ollama-setup.md}}
  \smallskip

  \small
  \begin{columns}[t]
    \begin{column}{.52\textwidth}
      \textbf{Components:}
      \begin{tabular}{ll}
        \toprule
        Component & Port \\
        \midrule
        OpenWebUI & \texttt{:3000} \\
        Ollama (native, Metal GPU) & \texttt{:11434} \\
        Model & \texttt{gemma4:e4b} (9.6\,GB) \\
        Native SPARQL tool & Workspace $\to$ Tools \\
        QLever & \texttt{:7030} \\
        SHMARQL & \texttt{:7032} \\
        MCPO (OpenWebUI only) & \texttt{:8001} \\
        \bottomrule
      \end{tabular}
    \end{column}
    \begin{column}{.46\textwidth}
      \textbf{Data flow:}
      \medskip

      \scriptsize
      \texttt{User} $\to$ \texttt{OpenWebUI} $\to$ \texttt{Ollama} \\
      \hspace{1em}$\to$ \textit{tool call} $\to$ \texttt{sparql\_query.py} \\
      \hspace{1em}$\to$ SPARQL GET $\to$ \texttt{QLever} \\[4pt]
      \texttt{User} $\to$ \texttt{SHMARQL} $\to$ \texttt{QLever} \\[4pt]
      \texttt{Claude Code} $\to$ MCP $\to$ \texttt{QLever}

    \end{column}
  \end{columns}
\end{frame}

% ---- 1.5 Deployment Options ----
\begin{frame}[shrink=5]{Deployment Options}
  \small
  \textbf{Problem:} Two QLever instances needed (DDB-EDM raw + mocho-aligned).
  At 27M objects, each instance requires $\geq$500\,GB storage;
  network traffic and public hosting costs are non-trivial.

  \smallskip
  \begin{columns}[t]
    \begin{column}{.48\textwidth}
      \textbf{Option A --- goethe-faust corpus only}
      \begin{itemize}\setlength{\itemsep}{2pt}
        \item Publish \textbf{115k-object} goethe-faust QLever + SHMARQL (already running locally)
        \item Offer full DDB corpus as \textbf{.nt download} (Zenodo / w3id)
        \item Reviewers can query live; full corpus available for bulk use
      \end{itemize}
      \smallskip
      \textit{Pros:} deployable today; low infra cost. \\
      \textit{Cons:} paper claims 27M --- live endpoint must match or abstract needs revision.
    \end{column}
    \begin{column}{.50\textwidth}
      \textbf{Option B --- full 27M on FIZ infra}
      \begin{itemize}\setlength{\itemsep}{2pt}
        \item Deploy on teach03 or FIZ server with public IP
        \item Both QLever instances exposed via SHMARQL + MCPO
        \item Persistent URI resolves to live SPARQL endpoint
      \end{itemize}
      \smallskip
      \textit{Pros:} matches paper claim; satisfies ISWC ``resource live at review'' requirement. \\
      \textit{Cons:} $\geq$1\,TB disk; infra ops burden; network traffic.
    \end{column}
  \end{columns}
  \begin{alertblock}{Open question}
    \small Can we get a public-facing FIZ server before 2 May?
    If not, Option A + Zenodo .nt is the fallback.
  \end{alertblock}
\end{frame}

% ---- 1.6 Critical Path Before 2 May ----
% source: roadmap-revised.md §8 Critical path + §9 Checklist
\begin{frame}{Critical Path --- Before 2 May (Abstract)}
  \small
  \begin{enumerate}\setlength{\itemsep}{5pt}
    \item \textbf{Register persistent URI} (w3id or Zenodo DOI) ---
          mandatory for Resource Availability Statement; \textbf{single most time-critical item}

    \item \textbf{QLever stats} ---
          triple/object count per named graph; provider count
          (\texttt{SELECT (COUNT(*) AS ?n) WHERE \{ ?s ?p ?o \}} per graph)

    \item \textbf{Document mocho PoC} ---
          subset size, classes produced, triple count, mapping quality notes

    \item \textbf{Adapt SHMARQL + \texttt{setup.sh}} from \texttt{goethe-faust/};
          point at GeMeA NT files; adjust ports and memory

    \item \textbf{Inventory MCP tools} exposed via \texttt{./setup.sh mcp-add};
          prepare 2 example agent interactions for §3.4

    \item \textbf{Draft abstract} with real numbers:
          scale + access methods + testbed framing + mocho PoC
  \end{enumerate}
  \begin{alertblock}{}
    \small SHMARQL screenshots, MCP examples, and §4 Quality prose can follow before 7 May ---
    but the resource must be live and core numbers real by 2 May.
  \end{alertblock}
\end{frame}


% =================================================================
\section{Mocho v1.1}
% source: mocho/notes/mocho-gatherer-plan.md
%         mocho/notes/mocho-gatherer-adr.md
%         mocho/notes/reasoner-change-plan.md
% =================================================================

% ---- 2.2 Context, current pattern, justification ----
% source: mocho-gatherer-adr.md §1 Context, §2 Justification, §3 Current pattern
\begin{frame}{DDB EDM $\to$ mocho: Context \& Current Pattern}
  \small
  \textbf{Context (§1):} sector + mediatype determine the \emph{domain} ontology;
  \texttt{htype} determines the \emph{domain class} within it. \\
  \texttt{dc:type} is uncontrolled (free-text literals) and varies per institution
  --- not alignable. Sector and mediatype carry centrally-defined vocnet IRIs:
  consistent semantics across all providers.

  \smallskip
  \textbf{Current pattern (v1.0 --- broken for DDB data):}
  \begin{center}
  \footnotesize
  \texttt{SubClassOf(mocho:Manifestation ObjectUnionOf(rdac:C10007 frbr:Manifestation \ldots))}
  \end{center}
  \vspace{-4pt}
  \textbf{Under the open-world assumption the reasoner cannot infer which disjunct holds.}
  DDB \texttt{edm:ProvidedCHO} objects carry no domain class types --- only \texttt{rdf:type edm:ProvidedCHO}.
  The union axiom contributes nothing to classification on raw EDM data.
\end{frame}

% ---- 2.2b Justification ----
% source: mocho-gatherer-adr.md §2 Justification v1.0 → v1.1
\begin{frame}{DDB EDM $\to$ mocho: Justification for v1.1}
  \small
  \textbf{Why switch from field-value dispatch to sector/mediatype-based restrictions? (§2)}
  \begin{itemize}\setlength{\itemsep}{4pt}
    \item Controlled fields (sector, mediatype) have the same semantics across all institutions
      $\Rightarrow$ safe to use as alignment signal
    \item \texttt{dc:type} fails this test: same field, different purposes per institution
    \item \texttt{htype} partially passes: centrally-defined codes map predictably to RiC-O and DoCO
    \item OWL \texttt{ObjectHasValue} restrictions on \texttt{mocho:sector}/\texttt{mocho:mediatype}
      serve as machine-readable specifications and SHACL validation rules
  \end{itemize}
\end{frame}

% ---- 2.2c Patterns ----
% source: mocho-gatherer-plan.md §1, Changes 5–7
\begin{frame}[shrink=5]{v1.1 Solution: Two OWL Patterns}
  \small
  % source: mocho-gatherer-plan.md summary table (Changes 5–7)
  \begin{tabular}{>{\bfseries}p{1.5cm}>{\raggedright\arraybackslash}p{5.5cm}>{\raggedright\arraybackslash}p{4.0cm}}
    \toprule
    & OWL construct & What it enables \\
    \midrule
    Pattern 1 &
      \texttt{SubClassOf(rdac:C10007 mocho:Manifestation)} \newline
      (one axiom per domain class, all WEMI levels) &
      \texttt{?x rdf:type mocho:Manifestation} returns \texttt{rdac:C10007},
      \texttt{aco:AudioManifestation}, \texttt{rico:Instantiation}, \ldots
      without listing each domain class explicitly \\[6pt]
    Pattern 2 &
      \texttt{SubObjectPropertyOf(} \newline
      \texttt{frbr:realization mocho:hasRealization)} \newline
      (domain props $\to$ mocho hub; Changes 5--6) &
      \texttt{?w mocho:hasRealization ?e} aggregates \texttt{frbr:realization},
      \texttt{rdaw:P10078}, \texttt{mo:performed\_in}, \ldots
      in a single cross-domain query \\
    \bottomrule
  \end{tabular}

  \smallskip
  \footnotesize
  \textbf{Pattern 1:} Any instance typed as a domain class is automatically inferred to be a mocho WEMI pivot. \\
  \textbf{Pattern 2:} mocho is the hub --- domain properties must be subproperties \emph{of} mocho properties;
  wrong direction causes inference to flow away from the hub.
\end{frame}

% ---- 2.2 Routing Table ----
% source: mocho-gatherer-adr.md §8 Reference: Sector x mediatype → domain class
\begin{frame}{Routing Table: Sector $\times$ Mediatype $\to$ Domain Class}
  \small
  % source: mocho-gatherer-adr.md §8
  \begin{center}
  \begin{tabular}{llll}
    \toprule
    Sector & Mediatype & Domain class & mocho pivot \\
    \midrule
    sparte002 Library  & any         & \texttt{rdac:C10007} ``Manifestation'' & \texttt{mocho:Manifestation} \\
    sparte006 Museum   & any         & \texttt{vra:Work}                      & \texttt{mocho:Work} / \texttt{mocho:Manifestation} \\
    sparte001 Archive  & any         & \texttt{rico:Instantiation}            & \texttt{mocho:Manifestation} \\
    sparte005 Media    & mt001 Audio & \texttt{aco:AudioManifestation}        & \texttt{mocho:Manifestation} \\
    sparte005 Media    & mt005 Video & \texttt{ec:MediaResource}              & \texttt{mocho:Manifestation} \\
    sparte006 Museum   & mt002 Photo & \texttt{vra:Image}                     & \texttt{mocho:Work} \\
    sparte002 Library  & mt002 Photo & \texttt{rdac:C10007}                   & \texttt{mocho:Manifestation} \\
    \bottomrule
  \end{tabular}
  \end{center}
  \smallskip
  \begin{block}{Note}
    \small Most DDB \texttt{edm:ProvidedCHO} objects land at \texttt{mocho:Manifestation} ---
    expected, as DDB holds physical/digital objects.
    Work-level mocho classes are populated separately via GND linking (Phase 0 pipeline).
  \end{block}
\end{frame}

% ---- 2.5b Changes Overview ----
% source: mocho-gatherer-plan.md §3 Summary of changes
\begin{frame}{Mocho v1.1 --- Changes Overview}
  \small
  \begin{tabular}{clp{6.5cm}ll}
    \toprule
    \# & & Change & Type & Count \\
    \midrule
    1 & & New prefixes (\texttt{vocnet-sparte:}, \texttt{vocnet-mtype:}, \texttt{vocnet-htype:})
        & ADD & 3 \\
    2 & & New \texttt{ObjectProperty mocho:sector} + domain/range/annotations
        & ADD & 1 prop \\
    3 & & New \texttt{ObjectProperty mocho:mediatype} + domain/range/annotations
        & ADD & 1 prop \\
    4 & & New declarations (2 OPs + 13 sparte/mt + 7 htype NamedIndividuals)
        & ADD & 22 \\
    5 & & \texttt{ClassAssertion} + bilingual labels for vocnet-htype individuals
        & ADD & 21 axioms \\
    6 & & \texttt{SubObjectPropertyOf} direction (FRBR + FaBiO) --- wrong $\to$ correct
        & REPLACE & 8 axioms \\
    7 & & New \texttt{SubObjectPropertyOf} (RDA, MO, RiC-O)
        & ADD & 4 axioms \\
    \bottomrule
  \end{tabular}
\end{frame}



% =================================================================
% =================================================================
\section{NER Annotation}
% source: rm057; gemea/notes/ner/
% =================================================================

% ---- 4.1 Status ----
% source: rm057 §5 Gold Dataset; §6 Training Dataset
\begin{frame}{NER Annotation --- Status}
  \small
  \begin{columns}[t]
    \begin{column}{.50\textwidth}
      \textbf{What was decided (RM 057):}
      \begin{itemize}\setlength{\itemsep}{2pt}
        \item Gold set: \textbf{395 records}, stratified by era $\times$ tier $\times$ \texttt{dc\_type}
        \item 4 strata: Pre-1700 / 1700--1800 / 19th-c / Modern
        \item Oversampled: LP (\textit{Leichenpredigt}) + EB (\textit{Einblattdruck})
        \item Labels: \texttt{TITLE}, \texttt{OTHER\_TITLE}, \texttt{PERSON}
        \item Path confirmed: LLM annotation + fine-tune \texttt{xlm-roberta-base}
      \end{itemize}
    \end{column}
    \begin{column}{.48\textwidth}
      \begin{alertblock}{Current status}
        Annotation \textbf{ongoing}. \\[3pt]
        Work \textbf{postponed} --- paper deadline (7 May) and mocho v1.1
        take priority. NER pipeline is \emph{not} in v1 paper scope;
        deferred to v1.1 (Phase 0).
      \end{alertblock}
    \end{column}
  \end{columns}
\end{frame}

% ---- 4.1b Corpus Reprocessing ----
% source: gemea/notes/project/reprocess-before-after.md
%         ADR-02 (pkl → parquet), ADR-03 (corpus migration)
\begin{frame}{NER Corpus --- Reprocessing: Before / After}
  \small
  % source: reprocess-before-after.md §1, §2, §3, §6
  \begin{center}
  \footnotesize
  \begin{tabular}{lll}
    \toprule
    & \textbf{Old (pkl)} & \textbf{New (parquet)} \\
    \midrule
    Rows          & 4.48M                          & \textbf{9.21M} \\
    Language filter & \texttt{dc:lang=ger} + \texttt{langid=ger} & \texttt{dc:lang} $\in$ \{ger, gmh, nds, lat\} \\
    htype filter  & notebook-defined               & ADR-01 BLANKET\_EXCLUDE \\
    Tokenizer     & spaCy (model unspecified)      & xlm-roberta-large BPE \\
    Median tokens & 8                              & 15 \\
    Reproducible  & No (notebook)                  & Yes (scripted) \\
    \midrule
    Silver tier-1 (training pool) & 336K (7.5\%)  & \textbf{689K} (7.5\%) \\
    \texttt{f\_other\_title} rate & 20.2\%         & \textbf{37.5\%} \\
    \texttt{has\_dot\_dash} rate  & 1.2\%          & 0.4\% \\
    \bottomrule
  \end{tabular}
  \end{center}
  \smallskip
  \begin{block}{What is unchanged}
    \footnotesize
    Tier proportions (0.1\,/\,7.5\,/\,92.4\%) stable --- ISBD heuristics scale uniformly. \quad
    SR-08 gold sample (395 records) present in new parquet --- \textbf{no re-sampling needed}.
  \end{block}
\end{frame}

% ---- 4.2 Done / Pending ----
% source: rm057 §5, §8 Next Steps
\begin{frame}{NER Annotation --- Done / Pending}
  \small
  \begin{columns}[t]
    \begin{column}{.48\textwidth}
      \textbf{Done:}
      \begin{itemize}\setlength{\itemsep}{2pt}
        \item[$\checkmark$] Gold sampling design (RM 057 §5)
        \item[$\checkmark$] Title boundary curation rules (\texttt{sr08})
        \item[$\checkmark$] Silver tier classification (3 tiers; $\sim$340K tier-1/2)
        \item[$\checkmark$] NuNER Zero ruled out (SR-09: F1\,=\,0.000)
        \item[$\checkmark$] Stratification rationale + CI sizing ($\pm$10\,pp)
      \end{itemize}
    \end{column}
    \begin{column}{.48\textwidth}
      \textbf{Pending (after 7 May):}
      \begin{itemize}\setlength{\itemsep}{2pt}
        \item[$\square$] LLM annotation of 395-record gold set
        \item[$\square$] Adjudication per \texttt{sr08\_title-boundary-curation.md}
        \item[$\square$] Fine-tune \texttt{xlm-roberta-base} on silver + LLM gold
        \item[$\square$] Evaluate: exact span F1, per label $\times$ per era
        \item[$\square$] Decision gate: pass $\geq$ targets $\Rightarrow$ integrate
              into \texttt{link\_gnd\_works.py}; fail $\Rightarrow$ error analysis
      \end{itemize}
    \end{column}
  \end{columns}
\end{frame}


% =================================================================
\section{Status \& Deliverables}
% source: gemea/paper/iswc-2026/
%         gemea/notes/project/roadmap-revised.md §5, §8, §9
% =================================================================

% ---- Paper Outline ----
% source: roadmap-revised.md §5 Paper Structure
\begin{frame}[shrink=10]{Paper Outline --- ISWC 2026 Resource Track ($\leq$15\,pp)}
  \footnotesize
  % source: roadmap-revised.md §5
  \begin{tabular}{clp{6.5cm}l}
    \toprule
    §\, & Section & Content & Status \\
    \midrule
    1   & Introduction  & Problem, gap, contributions, testbed framing       & \textcolor{orange}{draft} \\
    2   & Related Work  & Europeana LOD, ArtKB, BIBFRAME, comparison table   & \textcolor{orange}{mostly done} \\
    3   & Resource      & EDM schema, QLever stats, named graphs, VoID        & \textcolor{red}{needs stats} \\
    3.2 & Mocho PoC     & Alignment, subset size, class coverage, output     & \textcolor{red}{needs numbers} \\
    3.3 & SHMARQL       & Lightweight LD browser; Docker; screenshots         & \textcolor{red}{needs deploy} \\
    3.4 & MCP/MCPO      & Tools exposed; 2 example agent interactions         & \textcolor{red}{needs deploy} \\
    4   & Quality       & Triple counts, URI repair, mocho quality, limits    & \textcolor{red}{needs stats} \\
    5   & Usage         & 4 SPARQL examples; 2 MCP examples; use cases        & \textcolor{orange}{partial} \\
    6   & Impact        & DH + SW communities; testbed positioning            & \textcolor{green!60!black}{done} \\
    7   & Sustainability& Versioning, w3id URIs, open source, upstream dep.\  & \textcolor{green!60!black}{done} \\
    8   & Conclusion    & Summary, v2 directions                              & \textcolor{green!60!black}{draft} \\
        & Avail.\ Stmt  & Persistent URI + CC BY 4.0 + citation               & \textcolor{red}{needs URI} \\
    \bottomrule
  \end{tabular}
\end{frame}

% ---- Abstract ----
% source: gemea/paper/iswc-2026/00-abstract.tex
\begin{frame}{Abstract (due 2 May)}
  \small
  % source: gemea/paper/iswc-2026/00-abstract.tex
  \begin{block}{}
    We present GeMeA, a publicly accessible knowledge graph over 23M+ objects from the German
    Digital Library (DDB), queryable via SPARQL (QLever), browsable via SHMARQL, and accessible
    to AI agents through an MCP interface --- a combination not previously available for cultural
    heritage data at this scale. The resource includes a MOCHO-based ontology alignment layer
    mapping DDB EDM metadata to a domain- and modality-specific ontology, with partial GND entity
    linking via rule-based ISBD title extraction. GeMeA is released as a research testbed for work
    on ontology evaluation, agentic KG applications, RML-based declarative mapping, and provenance
    modelling of LLM-assisted enrichment. Data is licensed CC BY 4.0; the SPARQL endpoint, Linked
    Data browser, and source code are available at [URI].
  \end{block}
  \smallskip
  \begin{alertblock}{Before submitting}
    Replace \textbf{23M+} with verified count from QLever stats \quad$\cdot$\quad
    Replace \textbf{[URI]} with registered persistent URI
  \end{alertblock}
\end{frame}

% ---- Open TODOs ----
% source: roadmap-revised.md §8 Acceptance-critical items (ranked by weight)
\begin{frame}{Open TODOs --- Blocking Submission}
  \small
  % source: roadmap-revised.md §8 Acceptance-critical items (ranked by weight)
  \begin{enumerate}\setlength{\itemsep}{4pt}
    \item[$\square$] \textbf{Persistent URI registered} (w3id or Zenodo DOI) ---
          mandatory for Resource Availability Statement; \textbf{before 2 May}

    \item[$\square$] \textbf{Concrete stats in §3} --- triple counts, object counts,
          named graph inventory; every \texttt{\textbackslash todo\{\}} placeholder is a red flag

    \item[$\square$] \textbf{Mocho PoC quantified} --- subset size, classes produced,
          triple count; a PoC without numbers fails the quality criterion

    \item[$\square$] \textbf{SPARQL examples run against live endpoint} --- 4 queries
          in §5 must work; a SPARQL error poisons all subsequent reading

    \item[$\square$] \textbf{Resource live at review time} --- QLever + SHMARQL browser
          + persistent URI must all resolve; ISWC reviewers \emph{will} click the URI
  \end{enumerate}
\end{frame}

% ---- May 2 Checklist ----
% source: roadmap-revised.md §9 Before abstract (2 May)
\begin{frame}{Deliverables --- Before Abstract (2 May)}
  \small
  % source: roadmap-revised.md §9 Checklist / Before abstract (2 May)
  \begin{enumerate}\setlength{\itemsep}{5pt}
    \item[$\square$] \textbf{Persistent URI registered} (w3id or Zenodo DOI)
    \item[$\square$] \textbf{QLever stats}: triple/object count per named graph; provider count
    \item[$\square$] \textbf{Mocho PoC documented}: subset size, classes, triple count, quality
    \item[$\square$] \textbf{SHMARQL adapted + deployed}: copy from \texttt{goethe-faust/};
          adjust ports and memory; verify UI accessible
    \item[$\square$] \textbf{MCP tools inventory}: \texttt{./setup.sh mcp-add};
          list tools exposed; 2 example agent interactions drafted
    \item[$\square$] \textbf{SHMARQL screenshots} captured for §3.3 and §5
    \item[$\square$] \textbf{Abstract drafted} with real numbers:
          scale + access methods + testbed framing + mocho PoC
  \end{enumerate}
  \begin{alertblock}{}
    \small \textbf{Blocking all of the above:} persistent URI + live QLever endpoint.
  \end{alertblock}
\end{frame}

% ---- 5.2 May 7 Checklist ----
% source: roadmap-revised.md §9 Before full paper (7 May)
\begin{frame}[shrink=10]{Deliverables --- Before Full Paper (7 May)}
  \small
  % source: roadmap-revised.md §9 Before full paper (7 May)
  \begin{columns}[t]
    \begin{column}{.50\textwidth}
      \textbf{Paper sections:}
      \begin{itemize}\setlength{\itemsep}{2pt}
        \item[$\square$] §1 Introduction
        \item[$\square$] §2 Related Work --- finalize comparison table
        \item[$\square$] §3 Resource Description + VoID descriptor
        \item[$\square$] §3.2 Mocho PoC
        \item[$\square$] §3.3 SHMARQL --- deployment + screenshots
        \item[$\square$] §3.4 MCP/MCPO --- tools + 2 example interactions
        \item[$\square$] §4 Quality \& Validation
        \item[$\square$] §5 Usage --- 4 SPARQL + 2 MCP examples
        \item[$\checkmark$] §6 Impact
        \item[$\checkmark$] §7 Sustainability
        \item[$\checkmark$] §8 Conclusion (draft)
      \end{itemize}
    \end{column}
    \begin{column}{.48\textwidth}
      \textbf{Non-section deliverables:}
      \begin{itemize}\setlength{\itemsep}{2pt}
        \item[$\square$] Resource Availability Statement (URI + CC BY 4.0 + citation)
        \item[$\square$] Declaration of Use of Generative AI
        \item[$\square$] SPARQL examples tested on live endpoint
        \item[$\square$] Abstract finalized with verified numbers
      \end{itemize}
      \smallskip
      \textbf{Milestones:}
      \footnotesize
      \begin{tabular}{ll}
        \toprule
        Date & Event \\
        \midrule
        2 May   & Abstract \\
        7 May   & Full paper \\
        11--18 Jun & Rebuttal \\
        16 Jul  & Notification \\
        6 Aug   & Camera-ready \\
        27--29 Oct & ISWC 2026 \\
        \bottomrule
      \end{tabular}
    \end{column}
  \end{columns}
\end{frame}

% =================================================================
\end{document}
